\section{The complementary equation and its invariant graph}\label{sec:graph}

\subsection{The real elliptic operator}
Write $D_{0,p}=D'_p+d$ for the split Chern connection. On
trace-free Hermitian zero-forms set
\begin{equation}\label{eq:real-operator}
 L_p=\tfrac12(\Delta_{D'_p}+\Delta_d).
\end{equation}
Its real quadratic form is half the sum of the two derivative
energies. Its kernel therefore consists of parallel Hermitian
endomorphisms and equals $\ha$. Indeed, a parallel endomorphism
is holomorphic, the cross-Hom spaces between the distinct stable
factors vanish, and each stable factor has only scalar holomorphic
endomorphisms. Conversely, all multiplicity endomorphisms are
parallel. The restriction of $L_p$ to $Q_p$ has a common positive
gap for small $p$.

\begin{lemma}\label{lem:first-jet}
On the ambient coefficient ball the exact expression
\begin{equation}\label{eq:first-jet}
 \mathsf S_u^{-1}\M(p,x,u)
     =L_pu-b_p+\mathcal G_p(x,u,\nabla u)
\end{equation}
has a smooth remainder depending only on the first spatial jet
of $u$. For every $p$,
$\mathcal G_p(0,0)=0$ and
$D_{(x,u)}\mathcal G_p(0,0)=0$.
\end{lemma}

\begin{proof}
In \eqref{eq:chern}, second spatial derivatives occur only when
$e^{-u}D'_p(e^u)$ is differentiated. If $W$ denotes the
contracted second jet, its centered coefficient is
$e^{-u/2}(d\exp)_u(W)e^{-u/2}=\mathsf S_uW$.
Multiplying by $\mathsf S_u^{-1}$ cancels this coefficient
exactly. Every remaining derivative of $u$ has order at most one.
The finite form $a_p(x)$ and its spatial derivatives are smooth
and uniformly controlled on the chosen coefficient ball.

The constant term is $-b_p$. Put
$K_p=i\Lambda_pF_{D_{0,p}}$. Direct differentiation and the
zero-form K\"ahler identity give
\begin{align}
 D_u\M(p,0,0)u
 &=i\Lambda_p\ddbar D'_pu+\tfrac12[u,K_p]\notag\\
 &=\Delta_{D'_p}u-\tfrac12[K_p,u]
   =\tfrac12(\Delta_{D'_p}+\Delta_d)u. \label{eq:centered-linear}
\end{align}
Here $\Delta_d=\Delta_{D'_p}-[K_p,\cdot]$. The commutator
is needed when the factor slopes differ. Since
$\mathsf S_u^{-1}=I+O(u^2)$, it does not alter the
linearization, including in the presence of the constant $b_p$.
The coefficient derivative vanishes: contraction of
$D_{0,p}(x-x^*)$ is a combination of $d_p^*x$ and its adjoint,
which vanish for harmonic $x$. Subtracting the constant and
linear parts proves the statement.
\end{proof}

Substitution in \eqref{eq:slice} yields
\begin{equation}\label{eq:normal-system}
 \begin{split}
 \beta(p,x,u)&=R_p(u)^{-1}(b_p-P_p\mathcal G_p),\\
 \dot u&=-2L_pu+\N_p(x,u),\\
 \N_p&=-2\mathcal G_p+
       2(I-\mathsf B_uR_p(u)^{-1})b_p
       +2\mathsf B_uR_p(u)^{-1}P_p\mathcal G_p .
 \end{split}
\end{equation}
In the middle term $I$ includes $\ha$ in all sections.
The projection identity in Proposition~\ref{prop:slice}
shows that $\N_p$ takes values in $Q_p$.
Both $\mathsf B_u-I$ and $R_p(u)-I$ are quadratic, so
\begin{equation}\label{eq:vanishing-jets}
 D\N_p(0,0)=0,\qquad D_u\beta(p,0,0)=0.
\end{equation}
The finite linearization is $\Lambda_px=C_xb_p$, with norm
tending to zero as $p\to0$.

\subsection{A common parabolic domain}
Use \eqref{eq:complement-trivialization} to work on the fixed pair
\[
 \mathcal X_0=V_{\R}\oplus Q_0h^{3/8},\qquad
 \mathcal X_1=V_{\R}\oplus Q_0h^{19/8}.
\]
The continuous interpolation spaces needed below are
\begin{equation}\label{eq:interpolation}
 \begin{split}
 \mathcal X_\beta
  &=(\mathcal X_0,\mathcal X_1)^0_{5/8,\infty}
        =V_{\R}\oplus Q_0h^{13/8},\\
 \mathcal X_\alpha
  &=(\mathcal X_0,\mathcal X_1)^0_{7/8,\infty}
        =V_{\R}\oplus Q_0h^{17/8}.
 \end{split}
\end{equation}
We use the interpolation convention of
\cite[Proposition~2.2(c)]{shaosimonett}. One derivative of
$h^{13/8}$ belongs to the multiplication algebra $h^{5/8}$,
which embeds in $h^{3/8}$. Smooth multiplication and composition
therefore make the first-jet nonlinearities smooth
$\mathcal X_\beta\to\mathcal X_0$ maps. No restriction on
$\dim_{\R}X$ is required for this H\"older algebra estimate.
Smooth finite-rank projections preserve these interpolation spaces.
The exponents fix a convenient interpolation scale; no sharp
regularity threshold is asserted.

The transported operator
$\diag(0,2L_p|_{Q_p})$ on $\mathcal X_1\subset\mathcal X_0$
has continuous maximal regularity. We spell out the bundle
realization of the result used here. Complexify the real bundle
and choose a smooth complementary bundle making their sum trivial.
Extend $2L_p$ by a connection Laplacian on that complement with
the same scalar principal symbol. The resulting normally elliptic
system has domain $h^{19/8}$ on $h^{3/8}$ by the local resolvent
construction and Banach-valued extension of
\cite[Remark~3.3, Theorem~3.6 and Proposition~3.7]{shaosimonett}.
The original summand is reducing. Projection away from the
parallel kernel commutes with the operator, so its linear
solutions restrict to the required complement. A finite-dimensional
zero block also has maximal regularity. Uniform ellipticity and
smooth dependence on a compactly contained parameter neighborhood
give common constants.

In particular, on
$E_{0,p}=Q_ph^{3/8}$ and $E_{\alpha,p}=Q_ph^{17/8}$,
the semigroup of $A_p=2L_p|_{Q_p}$ satisfies
\begin{equation}\label{eq:semigroup}
 \|e^{-sA_p}\|_{E_\alpha\to E_\alpha}\le Me^{-\sigma s},
 \qquad
 \|e^{-sA_p}\|_{E_0\to E_\alpha}
       \le M(1+s^{-7/8})e^{-\sigma s}.
\end{equation}
Elliptic regularity identifies its spectrum with the smooth
$L^2$ spectrum. A common resolvent contour gives the exponential
bound from the fixed complementary gap; interpolation gives the
smoothing factor. The restricted semigroup is strongly continuous
on the little-H\"older initial space. The parameter-dependent
second-order term has remained in $L_p$ throughout this argument.

\subsection{The invariant graph}
\begin{proposition}\label{prop:graph}
For each prescribed finite $m\ge6$, after shrinking the parameter
and coefficient neighborhoods there is a jointly $C^m$,
$K$-equivariant map $\Psi(p,x)\in Q_ph^{19/8}$ whose graph
is locally invariant for \eqref{eq:slice}, and
\begin{equation}\label{eq:graph-bounds}
 \Psi(p,0)=D_x\Psi(p,0)=0,\qquad
 \|\Psi(p,x)\|_{h^{19/8}}\le C|x|^2,\qquad
 \|D_x\Psi(p,x)\|_{V\to h^{19/8}}\le C|x|.
\end{equation}
Its exact coefficient field is \eqref{eq:effective}, with
$\beta_0(p,x)=\beta(p,x,\Psi(p,x))$.
Every integrable trajectory, lifted by
$\dot g=\beta_0(p,x)g$, reconstructs an exact HYM solution.
\end{proposition}

\begin{proof}
At $p=0$, the finite space $V_{\R}$ is the entire center
space of the $(x,u)$ equation. The complementary linearization
$-2L_0|_{Q_0}$ is strictly stable. The equations
\eqref{eq:normal-system}--\eqref{eq:semigroup} verify the common
domain, smoothness, interpolation and continuous maximal-regularity
hypotheses of \cite[Theorem~2.1]{simonett1994}.
That parameter-dependent theorem gives a chosen $C^m$ graph
with values in the stable domain. The equilibrium $(x,u)=(0,0)$
exists for each parameter, and its construction gives
$\Psi(p,0)=0$.

Tangency for every nearby parameter requires an additional
argument. Set $T=D_x\Psi(p,0)$. Differentiating the exact
invariance equation, using \eqref{eq:vanishing-jets}, gives
\[
 T\Lambda_p=-2L_pT.
\]
The finite operator $\Lambda_p$ is diagonalizable with real
eigenvalues equal to differences of central block constants
in $b_p$. They tend to zero. Apply the displayed identity to
each eigenvector; the disjoint finite and normal spectra force
$T=0$. Taylor estimates on a smaller parameter product give
\eqref{eq:graph-bounds}.

Choose invariant norms and a $K$-invariant coefficient cutoff
in the fixed-point construction. Acting by $k\in K$ gives
another fixed point in the same uniqueness class, hence the
same graph. This proves equivariance without averaging
invariant graphs. Substitution gives the exact coefficient
field, and Proposition~\ref{prop:slice} gives reconstruction.
\end{proof}

The graph depends on the construction and on the chosen finite
regularity order. We use only its stated $C^m$ regularity.
Its invariance equation is
\begin{equation}\label{eq:invariance}
 D_x\Psi(p,x)C_x\beta_0(p,x)
          =-2L_p\Psi(p,x)+\N_p(x,\Psi(p,x)).
\end{equation}
The left side distinguishes it from a solution of a stationary
complementary curvature equation.

\subsection{The leading moment field}
\begin{proposition}\label{prop:moment}
The graph field satisfies \eqref{eq:effective-expansion},
with the moment convention \eqref{eq:moment}.
\end{proposition}

\begin{proof}
Transfer the connection of $h_pe^\Psi$ isometrically to $h_p$.
Its $(0,1)$ perturbation has the expansion
\[
 \widetilde a=x+a_2(x)-\tfrac12d\Psi_2(p,x)+O(|x|^3),
 \qquad a_2(x)=-d_p^*G_{2,p}(x\wedge x).
\]
The quadratic correction is orthogonal to every harmonic
one-form: its two terms are coexact and exact, respectively.
For a parallel Hermitian endomorphism $\xi$, integration of
curvature gives the exact identity
\begin{equation}\label{eq:projected-curvature}
 \langle P_p\M(p,x,\Psi),\xi\rangle
 =-\langle b_p,\xi\rangle+
          \operatorname{Re}\langle[\xi,\widetilde a],\widetilde a\rangle.
\end{equation}
The derivative term integrates to zero because $\xi$ is
parallel. In a unitary coframe the remaining term is a sum
of $[A_j,A_j^*]$, and
$\tr(\xi[A_j,A_j^*])
 =\operatorname{Re}\tr([\xi,A_j]A_j^*)$.
This also verifies the sign and requires no integrability
assumption on the ambient coefficient.

The action $[\xi,\cdot]$ preserves harmonic forms and is
self-adjoint for the real one-form inner product. Consequently
both cubic cross terms between $x$ and the quadratic correction
vanish. Equation~\eqref{eq:projected-curvature} becomes
$P_p\M=-b_p+\mu_p^\sharp+O(|x|^4)$.
Moreover $\mathsf S_\Psi^{-1}-I$ and $R_p(\Psi)^{-1}-I$
are $O(|x|^4)$, whereas
$\M(p,x,\Psi)=O(|p|+|x|^2)$ in the required spatial norm.
The formula for $\beta$ in \eqref{eq:slice} therefore retains
the quartic remainder. The vanished jets and the $C^m$ bounds
give its first two differentiated estimates.

Finally,
$d\langle\mu_p^\sharp,\xi\rangle[v]
 =2\operatorname{Re}\langle C_x\xi,v\rangle$.
Differentiating $\frac14\|\mu_p^\sharp-b_p\|^2$ gives
the leading field in Theorem~\ref{thm:main}.
\end{proof}

The factor-rank and volume weights are part of the Riesz
identification in \eqref{eq:moment}. Changing the moment
normalization changes the scalar time factor. In all
tracking statements we use the exact field $\beta_0$,
including its remainder.
